% Proposed B.111. Exact q0 witness; no full annular centralizer conclusion.
\statement{B.111}{Proved}{Restoring the retained zero loop kills a nontrivial diagonal in the actual slit-region group}
\label{stmt:B111}
Let $R_{\mathrm{slit}}$ be the actual inner bad-region orbifold group
over the simply connected slit base of B.110, and let
$R_{\mathrm{ann}}$ be the corresponding group over $0<|q|<1$.
Use the based inclusion and the same fiber marking throughout.
In the HNN presentation of B.110 write
\[
 \overline{x}=cxc,\qquad c^2=1,\qquad
 w p w^{-1}=e,\qquad w r w^{-1}=\overline b.
\]
Then the inclusion-induced homomorphism
\[
 j:R_{\mathrm{slit}}\longrightarrow R_{\mathrm{ann}}
\]
kills the nontrivial element $[p,\overline p]$.
In particular, $j$ is not injective.  The same element is nontrivial
in $\widehat{R}_{\mathrm{slit}}$ and is killed by $\widehat j$.
This proves neither a complete simplified presentation of
$R_{\mathrm{ann}}$ nor an equality of its cone centralizer with its
source image.

\begin{proof}
\emph{B.111.1. The actual retained loop and the chosen contact.}
The initial six-meridian order in B.110 is
\[
 (x_1,x_2,x_3,x_4,x_5,x_6)
 =(h_1^2,\text{cone }0,h_4^2,h_2^2,h_3^2,q),
 \qquad x_2=c.
\]
Here the right-hand side names the corresponding punctures; the
$x_i$ are their based positive meridians.
The actual positive $q=0$ lasso has chronological braid
$W_0\,(3,3)\,W_0^{-1}$, where
\begin{equation}\label{eq:B111-whisker}
 W_0=(2,1,-2,-2,4,5,-4).
\end{equation}
For the calibrated Artin maps
\[
 \phi_i(x_i)=x_i x_{i+1}x_i^{-1},\qquad
 \phi_i(x_{i+1})=x_i,
\]
with the other generators fixed, chronological push means that the
first crossing acts first.  Thus
\begin{equation}\label{eq:B111-q0-action}
 \psi_0=\operatorname{back}(W_0)\circ\phi_3^2
                    \circ\operatorname{push}(W_0).
\end{equation}
The common far-section generator $g_0$ satisfies
$g_0^{-1}xg_0=\psi_0(x)$.
The exact path, crossing and deformation certificates of B.110
apply to \eqref{eq:B111-whisker}; the loop is retained, not filled.

Put
\begin{equation}\label{eq:B111-raw-contact}
 U=x_3^{-1}x_1c x_1^{-1}x_3,
 \qquad r_{33}=[x_5,Ux_5U],
 \qquad V=x_5x_6^{-1}x_5^{-1}.
\end{equation}
The element $U$ is an involution.  The relation $r_{33}=1$ is the
actual fourth-twist contact at $\beta_{33}$, in the same whisker
marking, not an independently conjugated peripheral relation.
Direct substitution in \eqref{eq:B111-q0-action} gives the two short
identities
\begin{equation}\label{eq:B111-inverse-images}
 \psi_0^{-1}(x_5)=x_5,
 \qquad \psi_0^{-1}(U)=VUV^{-1}.
\end{equation}
These identities are already identities in $C_2*F_5$; they use no
filled-node relation.

\emph{B.111.2. Literal conversion to the HNN marking.}
Use the same-sheet Nielsen basis of B.110.  The entries needed here
are
\[
 a=x_1,\qquad b=x_1^{-1}x_3x_1,\qquad
 e=x_1^{-1}x_3x_5x_3^{-1}x_1,\qquad f=x_6.
\]
Set
\[
 g=ab^{-1},\qquad S=acac=a\overline a.
\]
The already verified inverse changes of generators give
\begin{equation}\label{eq:B111-marked-substitutions}
 x_5=g e g^{-1},\qquad U=g c g^{-1},\qquad
 f=g e^{-1}S^{-1}w^{-1}g^{-1}.
\end{equation}
The base relation $[a,\overline a]=1$ implies
$cSc=\overline a a=S$, so $S$ commutes with $c$.
In particular, \eqref{eq:B111-marked-substitutions} implies
\[
 V=g e w S g^{-1},\qquad
 \psi_0^{-1}(U)=g e w c w^{-1}e^{-1}g^{-1}.
\]
Writing $Q=wcw^{-1}$ and using $Q^2=1$, we obtain
\begin{align}
 \psi_0^{-1}(r_{33})
   &=g[e,eQeQe^{-1}]g^{-1}\\
   &=(g e)[e,QeQ](g e)^{-1}\\
   &=(g e w)[p,\overline p](g e w)^{-1}\\
   &=(a b^{-1}w p)[p,\overline p]
                         (a b^{-1}w p)^{-1}.
 \label{eq:B111-exact-witness}
\end{align}
The third line uses $e=wpw^{-1}$ and hence
$QeQ=w\overline p w^{-1}$; the last uses $ew=wp$.
This derivation retains the full conjugating word and uses only the
displayed base commutator and HNN identification.

\emph{B.111.3. Nontriviality before restoring the loop.}
The two vertices $p,\overline p$ of the twelve-vertex graph in B.110
are not joined by an edge.  Killing every other graph vertex defines
a coordinate retraction
\[
 \Gamma\longrightarrow F(P,Q),\qquad
 p\longmapsto P,\quad\overline p\longmapsto Q.
\]
It sends $[p,\overline p]$ to the reduced nonempty word
$PQP^{-1}Q^{-1}$.  Thus this element is nontrivial in $\Gamma$,
in $\Gamma\rtimes\langle c\rangle$, and, by the actual HNN
base injection proved in B.110, in $R_{\mathrm{slit}}$.
The residual finiteness of $\Gamma$ and the full induced topology
on the HNN base proved in B.110 also show that it is nontrivial in
$\widehat{R}_{\mathrm{slit}}$.

\emph{B.111.4. The actual annular relation.}
Van Kampen over the punctured disk retains $g_0$ and imposes
\[
 g_0^{-1}x_i g_0=\psi_0(x_i)\quad(1\leq i\leq6),
\]
in addition to the ten filled relations of the slit region and $c^2=1$.
Consequently
\[
 \psi_0^{-1}(r_{33})=g_0r_{33}g_0^{-1}=1
                 \quad\text{in }R_{\mathrm{ann}}.
\]
Equation~\eqref{eq:B111-exact-witness} therefore proves
$j([p,\overline p])=1$.  This equality persists under full profinite
completion.  No injection of a post-filling fiber group, stability of
the ten-relation normal closure under $\psi_0$, or additional annular
normal-form assertion has been used.

\emph{B.111.5. Exact records.}
The separate literal certificate checks both inverse monodromy arrays
on all six generators, the two identities
\eqref{eq:B111-inverse-images}, the marked substitutions, the complete
conjugacy \eqref{eq:B111-exact-witness}, and the coordinate retraction.
The earlier twenty-image comparison is preserved separately.
The genuine records are
\begin{raggedright}\small
\path{work_CE/code/certify_inner_q0_diagonal22.py},
its \path{.out} and \path{.err}, and
\path{work_CE/next_symmetric_nodes/inner_q0_diagonal22.json};
the earlier producer is
\path{work_CE/code/check_inner_q0_HNN22.py}, with its output and
\path{work_CE/next_symmetric_nodes/inner_q0_HNN22.json}.
\par\end{raggedright}
\end{proof}
